\section{Tietze 扩张定理}\label{section:ch04-tietze-extension}

Tietze 扩张定理是 Urysohn 引理的一个直接推论，这是一个很有用的定理。它所处理的问题是，将定义在 $X$ 的一个子空间上的实值连续函数扩张成为定义在整个空间 $X$ 上的连续函数。这个定理在拓扑学的许多应用中很重要。

\begin{theorem}[Tietze 扩张定理]\label{theorem:ch04-tietze-extension}
设 $X$ 是一个正规空间。$A$ 是 $X$ 的一个闭子集。则
\begin{enumerate}
\item 任何一个从 $A$ 到 $\mathbb{R}$ 的闭区间 $[a,b]$ 中的连续映射都可以扩张为从整个空间 $X$ 到 $[a,b]$ 中的一个连续映射.
\item 任何一个从 $A$ 到 $\mathbb{R}$ 中的连续映射都可以扩张为从整个空间 $X$ 到 $\mathbb{R}$ 中的一个连续映射.
\end{enumerate}
\end{theorem}

\par
\begin{figure}[htbp]
\centering
\includegraphics[width=0.6\textwidth]{ch04-fig35-1.png}
\caption{Tietze 扩张定理中的逼近函数 $g$}\label{figure:ch04-tietze-approximation}
\end{figure}
\begin{proof}
定理证明的思路是构造一个定义在整个空间 $X$ 上的连续函数序列 $s_{n}$，使得序列 $s_{n}$ 一致收敛，并且 $s_{n}$ 在 $A$ 上的限制随着 $n$ 的增大愈来愈逼近 $f$。于是其极限函数是连续的，并且它在 $A$ 上的限制等于 $f$。

第一步。这一步是在整个 $X$ 上构造一个特别的函数 $g$，使得 $g$ 不太大并且在集合 $A$ 上可以相当精确地逼近 $f$。更确切地说，我们考虑 $f\colon A\to[-r,r]$ 的情形。我们断言：存在一个连续函数 $g\colon X\to\mathbb{R}$ 使得
\begin{align*}
|g(x)|\leqslant\frac{1}{3}r\quad\text{对于所有 }x\in X,
\end{align*}
\begin{align*}
|g(a)-f(a)|\leqslant\frac{2}{3}r\quad\text{对于所有 }a\in A.
\end{align*}
函数 $g$ 的构造如下：

将区间 $[-r,r]$ 分成三个长度都等于 $\frac{2}{3}r$ 的区间。
\begin{align*}
I_{1}=\big[-r,-\frac{1}{3}r\big],\quad I_{2}=\big[-\frac{1}{3}r,\frac{1}{3}r\big],\quad I_{3}=\big[\frac{1}{3}r,r\big].
\end{align*}
令 $A$ 的子集 $B$ 和 $C$ 分别为
\begin{align*}
B=f^{-1}(I_{1})\quad\text{和}\quad C=f^{-1}(I_{3}).
\end{align*}
因为 $f$ 是连续的，所以 $B$ 和 $C$ 是 $A$ 中无交的闭子集。因此，它们是 $X$ 的闭集。根据 Urysohn 引理，存在一个连续函数
\begin{align*}
g\colon X\longrightarrow\big[-\frac{1}{3}r,\frac{1}{3}r\big]
\end{align*}
满足条件：对于 $B$ 中每一点 $x$，$g(x)=-\frac{1}{3}r$，对于 $C$ 中每一点 $x$，$g(x)=\frac{1}{3}r$。

于是对于所有的 $x$，$|g(x)|\leqslant\frac{1}{3}r$。我们断言：对于 $A$ 中的每一个点 $a$，
\begin{align*}
|g(a)-f(a)|\leqslant\frac{2}{3}r.
\end{align*}
这里有三种情况。若 $a\in B$，则 $f(a)$ 和 $g(a)$ 都在 $I_{1}$ 中。若 $a\in C$，则 $f(a)$ 和 $g(a)$ 都在 $I_{3}$ 中。若 $a\notin B\cup C$，则 $f(a)$ 和 $g(a)$ 都在 $I_{2}$ 中。因此，无论哪种情形，都有 $|g(a)-f(a)|\leqslant\frac{2}{3}r$。参见\cref{figure:ch04-tietze-approximation}。

第二步。现在我们来证明 Tietze 定理中的结论 (a)。不失一般性，我们可以将 $\mathbb{R}$ 的任意闭区间 $[a,b]$ 改换成区间 $[-1,1]$。

设 $f\colon X\to[-1,1]$ 是一个连续映射。则 $f$ 满足第一步中的假设的要求，其中 $r=1$。因此，存在一个定义在 $X$ 上的连续实值函数 $g_{1}$，使得
\begin{align*}
|g_{1}(x)|\leqslant 1/3\quad\text{对于 }x\in X,
\end{align*}
\begin{align*}
|f(a)-g_{1}(a)|\leqslant 2/3\quad\text{对于 }a\in A.
\end{align*}
考虑函数 $f-g_{1}$。它是从 $A$ 到 $[-2/3,2/3]$ 中的映射，从而我们可以对于 $r=2/3$ 再次应用第一步的结果。我们得到一个定义在整个空间 $X$ 上的实值函数 $g_{2}$，使得
\begin{align*}
|g_{2}(x)|\leqslant\frac{1}{3}\Big(\frac{2}{3}\Big)\quad\text{对于 }x\in X,
\end{align*}
\begin{align*}
|f(a)-g_{1}(a)-g_{2}(a)|\leqslant\Big(\frac{2}{3}\Big)^{2}\quad\text{对于 }a\in A.
\end{align*}
然后，我们再对函数 $f-g_{1}-g_{2}$ 应用第一步的结果。依此类推。

一般地，我们有定义在整个空间 $X$ 上的实值函数 $g_{1},\cdots,g_{n}$，使得对于 $a\in A$ 有
\begin{align*}
|f(a)-g_{1}(a)-\cdots-g_{n}(a)|\leqslant\Big(\frac{2}{3}\Big)^{n}.
\end{align*}
对于函数 $f-g_{1}-\cdots-g_{n}$ 以及 $r=\big(\frac{2}{3}\big)^{n}$ 应用第一步中的结论，我们得到一个定义在整个空间 $X$ 上的实值函数 $g_{n+1}$，使得
\begin{align*}
|g_{n+1}(x)|\leqslant\frac{1}{3}\Big(\frac{2}{3}\Big)^{n}\quad\text{对于 }x\in X,
\end{align*}
\begin{align*}
|f(a)-g_{1}(a)-\cdots-g_{n+1}(a)|\leqslant\Big(\frac{2}{3}\Big)^{n+1}\quad\text{对于 }a\in A.
\end{align*}
根据归纳原则，对于每一个 $n$ 定义了 $g_{n}$。

对于 $X$ 中的每一点 $x$，我们现在定义
\begin{align*}
g(x)=\sum_{n=1}^{\infty}g_{n}(x).
\end{align*}
当然，我们还必须知道这个无穷级数的收敛性。而我们可以通过将其与几何级数
\begin{align*}
\frac{1}{3}\sum_{n=1}^{\infty}\Big(\frac{2}{3}\Big)^{n-1}
\end{align*}
进行比较，应用微积分中的比较定理，得知这个级数是收敛的。

为了证明 $g$ 的连续性，我们需要证明序列 $s_{n}$ 一致收敛到 $g$。根据数学分析中的``Weierstrass M-判别法''便可得到这个序列的一致收敛性。如果不用这个判别法，我们只须注意：当 $k>n$ 时有
\begin{align*}
|s_{k}(x)-s_{n}(x)|=\Big|\sum_{i=n+1}^{k}g_{i}(x)\Big|
\end{align*}
\begin{align*}
\leqslant\frac{1}{3}\sum_{i=n+1}^{k}\Big(\frac{2}{3}\Big)^{i-1}
\end{align*}
\begin{align*}
<\frac{1}{3}\sum_{i=n+1}^{\infty}\Big(\frac{2}{3}\Big)^{i-1}=\Big(\frac{2}{3}\Big)^{n}.
\end{align*}
固定 $n$，并且令 $k\to\infty$，可见对于任何一个 $x\in X$ 有
\begin{align*}
|g(x)-s_{n}(x)|\leqslant\Big(\frac{2}{3}\Big)^{n}.
\end{align*}
因此 $s_{n}$ 一致收敛到 $g$。

以下证明：对于 $a\in A$，$g(a)=f(a)$。设 $s_{n}(x)=\sum_{i=1}^{n}g_{i}(x)$，即级数前 $n$ 项的部分和。那么 $g(x)$ 是部分和的无穷序列 $s_{n}(x)$ 的极限。由于对于所有 $A$ 中的点 $a$，有
\begin{align*}
\Big|f(a)-\sum_{i=1}^{n}g_{i}(a)\Big|=|f(a)-s_{n}(a)|\leqslant\Big(\frac{2}{3}\Big)^{n},
\end{align*}
所以对于所有的 $a\in A$ 有 $s_{n}(a)\to f(a)$。从而，对于 $a\in A$ 有 $f(a)=g(a)$。

最后，我们来证明 $g$ 是从 $X$ 到 $[-1,1]$ 中的映射。事实上，它可以由级数 $(1/3)\sum(2/3)^{n}$ 收敛于 $1$ 而轻松获得。然而，它不过是证明中的一个幸运巧合，而非证明的实质部分。倘若，我们已知的仅是 $g$ 为从 $X$ 到 $\mathbb{R}$ 中的一个映射，映射 $r\colon\mathbb{R}\to[-1,1]$ 为
\begin{align*}
r(y)=y\quad\text{如果 }|y|\leqslant 1,
\end{align*}
\begin{align*}
r(y)=y/|y|\quad\text{如果 }|y|\geqslant 1,
\end{align*}
则从 $X$ 到 $[-1,1]$ 的映射 $r\circ g$ 便是 $f$ 的一个扩张。

第三步。现在我们来证明定理中的结论 (b)，这时 $f$ 是从 $A$ 到 $\mathbb{R}$ 的一个映射。我们可以用 $(-1,1)$ 来代替 $\mathbb{R}$，因为这个区间同胚于 $\mathbb{R}$。

因而，令 $f$ 是从 $A$ 到 $(-1,1)$ 的一个连续映射。已经证明了的 Tietze 定理的前一部分告诉我们：$f$ 可以扩张为从 $X$ 到闭区间的一个连续映射 $g\colon X\to[-1,1]$。如何找到一个映射 $h$ 将 $X$ 映到这个开区间呢？

对于给定的 $g$，定义一个 $X$ 的子集 $D$
\begin{align*}
D=g^{-1}(\{-1\})\cup g^{-1}(\{1\}).
\end{align*}
根据 $g$ 的连续性可见，$D$ 是 $X$ 的一个闭子集。由于 $g(A)=f(A)$ 包含于 $(-1,1)$，所以集合 $A$ 与 $D$ 无交。再根据 Urysohn 引理，存在一个连续函数 $\phi\colon X\to[0,1]$ 使得 $\phi(D)=\{0\}$ 和 $\phi(A)=\{1\}$。定义
\begin{align*}
h(x)=\phi(x)g(x).
\end{align*}
$h$ 作为两个连续映射的乘积是连续的。由于对于 $A$ 中的每一个点 $a$ 有
\begin{align*}
h(a)=\phi(a)g(a)=1\cdot g(a)=f(a),
\end{align*}
所以 $h$ 也是 $f$ 的一个扩张。最后，$h$ 将整个空间 $X$ 映入开区间 $(-1,1)$ 中。因为若 $x\in D$，则 $h(x)=0\cdot g(x)=0$；而当 $x\notin D$ 时，$|g(x)|<1$。由此可见 $|h(x)|\leqslant 1\cdot|g(x)|<1$。
\end{proof}

\subsection*{习题}

\begin{enumerate}
\item 证明 Tietze 扩张定理蕴涵着 Urysohn 引理。
\item 在 Tietze 扩张定理的证明中，第一步将区间 $[-r,r]$ 三等分，这一巧妙方法的实质是什么？倘若将以上方式改为把这个区间拆分成以下三个区间
\begin{align*}
I_{1}=[-r,-ar],\quad I_{2}=[-ar,ar],\quad I_{3}=[ar,r],
\end{align*}
其中 $a$ 是满足 $0<a<1$ 的某一个实数。试问除 $a=1/3$ 以外，$a$ 还有哪些值（如果存在的话）可以用于完成定理的证明？
\item 设 $X$ 是可度量化的。证明以下命题等价：
\begin{enumerate}
\item $X$ 相对于每一个导出 $X$ 的拓扑的度量是有界的.
\item 每一连续函数 $\phi\colon X\to\mathbb{R}$ 是有界的.
\item $X$ 是极限点紧致的.
\end{enumerate}
[提示：若 $\phi\colon X\to\mathbb{R}$ 是一个连续函数，则 $F(x)=x\times\phi(x)$ 是从 $X$ 到 $X\times\mathbb{R}$ 中的一个嵌入。如果 $A$ 是 $X$ 的一个没有极限点的无限子集，则 $\phi$ 是从 $A$ 到 $\mathbb{Z}_{+}$ 的一个连续满射。]
\item 设 $Z$ 是一个拓扑空间。若 $Y$ 是 $Z$ 的子空间，我们说 $Y$ 是 $Z$ 的一个\textbf{收缩核}（retract），如果存在一个连续映射 $r\colon Z\to Y$ 使得对于任何一个 $y\in Y$ 有 $r(y)=y$。
\begin{enumerate}
\item 证明：若 $Z$ 是一个 Hausdorff 空间，并且 $Y$ 是 $Z$ 的一个收缩核，则 $Y$ 是 $Z$ 的一个闭子集.
\item 设 $A$ 是由 $\mathbb{R}^{2}$ 中两个点组成的子集，证明 $A$ 不是 $\mathbb{R}^{2}$ 的收缩核.
\item 设 $S^{1}$ 为 $\mathbb{R}^{2}$ 中的单位圆周。证明：$S^{1}$ 是 $\mathbb{R}^{2}-\{\boldsymbol{0}\}$ 的一个收缩核，其中 $\boldsymbol{0}$ 是原点。请判断 $S^{1}$ 是否为 $\mathbb{R}^{2}$ 的收缩核？
\end{enumerate}
\item 一个空间 $Y$ 称为具有\textbf{通有扩张性质}（universal extension property），若对于任意给定的由正规空间 $X$、$X$ 的闭子集 $A$ 以及一个连续函数 $f\colon A\to Y$ 所组成的每一个三元组，存在一个从 $X$ 到 $Y$ 的连续映射为 $f$ 的扩充。
\begin{enumerate}
\item 证明 $\mathbb{R}^{J}$ 具有通有扩张性质.
\item 证明：若 $Y$ 同胚于 $\mathbb{R}^{J}$ 的一个收缩核，则 $Y$ 也具有通有扩张性质.
\end{enumerate}
\item 设 $Y$ 是一个正规空间。称 $Y$ 是一个\textbf{绝对收缩核}（absolute retract），是指对于每一个空间偶对 $(Y_{0},Z)$，若 $Z$ 是一个正规空间并且 $Y_{0}$ 是 $Z$ 中同胚于 $Y$ 的一个闭子空间，则空间 $Y_{0}$ 是 $Z$ 的一个收缩核。
\begin{enumerate}
\item 证明：若 $Y$ 具有通有扩张性质，则 $Y$ 是一个绝对收缩核.
\item 证明：若 $Y$ 是一个绝对收缩核并且 $Y$ 是紧致的，则 $Y$ 具有通有扩张性质。[提示：假设 Tychonoff 定理成立，从而有 $[0,1]^{J}$ 是正规的。将 $Y$ 嵌入到 $[0,1]^{J}$。]
\end{enumerate}
\item (a) 证明对数螺线
\begin{align*}
C=\{\,0\times 0\,\}\cup\{\,e^{t}\cos t\times e^{t}\sin t\mid t\in\mathbb{R}\,\}
\end{align*}
是 $\mathbb{R}^{2}$ 的一个收缩核。你能给出这样的一个收缩 $r\colon\mathbb{R}^{2}\to C$ 吗？

(b) 证明由图 35.2 所示的``$x$-轴扭结'' $K$ 是 $\mathbb{R}^{3}$ 的一个收缩。

\textbf{定理}\quad 设 $Y$ 是一个正规空间。则 $Y$ 是一个绝对收缩核当且仅当 $Y$ 有通有扩张性质。

[提示：若 $X$ 和 $Y$ 是无交的两个正规空间，$A$ 是 $X$ 的一个闭子集，并且 $f\colon A\to Y$ 是一个连续映射，则\textbf{贴附空间}（adjunction space）$Z_{f}$ 被定义为一个商空间，它是在 $X\cup Y$ 中黏合 $A$ 中点 $a$、$f(a)$ 和 $f^{-1}(\{f(a)\})$ 的所有点所得到的空间。应用 Tietze 扩张定理，证明 $Z_{f}$ 是正规的。若 $p\colon X\cup Y\to Z_{f}$ 是一个商映射，证明 $p\mid Y$ 是 $Y$ 与 $Z_{f}$ 的一个闭子空间之间的一个同胚。]
\item 设 $X_{1}\subset X_{2}\subset\cdots$ 是空间的一个序列，其中，对于每一个 $i$，$X_{i}$ 是 $X_{i+1}$ 的一个闭子空间。设 $X$ 是所有 $X_{i}$ 的并。我们可以将 $X$ 拓扑化：$X$ 的子集 $U$ 称为 $X$ 的一个开集，如果对于每一个 $i$，$U\cap X_{i}$ 在 $X_{i}$ 中是开的。
\begin{enumerate}
\item 证明：如上定义的开集确实决定了 $X$ 上的一个拓扑，并且每一个空间 $X_{i}$ 关于这一拓扑是 $X$ 的子空间（事实上，是闭子空间）。这个拓扑称为与各个子空间 $X_{i}$ 相通的（coherent）拓扑.
\item 证明：若对于每一个 $i$，$f\mid X_{i}$ 连续，则 $f\colon X\to Y$ 也连续.
\item 证明：若每一个空间 $X_{i}$ 都是正规的，则 $X$ 也是正规的。[提示：任意给定 $X$ 的两个无交的闭子集 $A$ 和 $B$，令 $f$ 在 $A$ 上等于 $0$，在 $B$ 上等于 $1$，然后将 $f$ 依次扩张到 $A\cup B\cup X_{i}$ 上，其中 $i=1,2,\cdots$。]
\end{enumerate}
\end{enumerate}